\documentclass[10pt,a4paper]{amsart}
\usepackage[margin=19mm]{geometry}
\usepackage{amsmath,amssymb,mathtools}
\usepackage[scr=rsfs,cal=euler]{mathalfa}
\usepackage{xfrac}
\usepackage[unicode,hidelinks,bookmarks=false]{hyperref}
\newcommand{\ie}{i.e.\ }
\newcommand{\cf}{cf.\ }
\newcommand{\resp}{respectively\ }
\newcommand{\Subsection}{\subsection}
\providecommand{\nobreakdash}{\nobreak\hbox{-}}
\setlength{\emergencystretch}{2em}
\multlinegap0pt
\usepackage[shortlabels]{enumitem}
\usepackage{amssymb}
\usepackage{mathtools}
\newtheorem{lemma}{Lemma}
\newtheorem{theorem}{Theorem}
\newcommand{\Id}{\mathrm{Id}}
\newcommand{\K}{\mathbb K}
\newcommand{\N}{\mathbb N}
\title{Finite-codimensional subspaces of Daugavet spaces: projection constants and minimal projections}
\author{Tomasz Kania, Grzegorz Lewicki}
\date{Source: arXiv:2604.10771v1 (CC BY 4.0)}
\begin{document}
\maketitle

\noindent\emph{Source and licence.} Finite-codimensional subspaces of Daugavet spaces: projection constants and minimal projections. Version arXiv:2604.10771v1 (CC BY 4.0), distributed by arXiv under the Creative Commons Attribution 4.0 International licence, http://creativecommons.org/licenses/by/4.0/, as recorded in the arXiv licence metadata for this exact version and shown on the abstract page https://arxiv.org/abs/2604.10771v1. Full licence notice: \texttt{licenses/arxiv-2604.10771-CC-BY-4.0.txt}.\par\smallskip

\noindent\textit{Editorial scope.} Counted unit: the article's theorem thm:piecewise-main (source0.tex lines 608--620). The article's complete original proof of it is delivered with it (source0.tex lines 622--737, one \texttt{proof} environment of 116 lines). No mathematics is written, repaired or re-derived: the statement and the proof are the article's own lines, in the article's own notation, and no retained line is substituted or re-flowed.  Retained uncounted as the source-local context the proof needs: lines 247--254; lines 594--597.  Nothing was omitted from any retained span, so the package carries no omission record.  No retained line contains a citation, so the wrapper carries no bibliography and no bibliography entry.  No retained line contains a figure environment, an image-inclusion command or drawing code, so no image is delivered and the package declares no asset dependency.  Editorial formatting only, in addition: an AMS \texttt{amsart} preamble that restates the article's own declarations and macros used by the retained lines, namely the environments \texttt{lemma}, \texttt{theorem} on separate counters, together with the standard packages those lines need.  The retained environments print in this document as Lemma 1, Theorem 1 in order of appearance, whereas the article prints the same items with the numbering of its own full document.  The complete native source of the article as distributed in the arXiv e-print for this exact version is bundled unchanged as \texttt{sources/198291/source0.tex}.
\begin{lemma}
\label{lem:wstar-coded}
Let $K$ be a compact Hausdorff space, let $M(K)=C(K)^*$, and let $E\subset M(K)$ be finite-dimensional. A finite-rank operator $P\colon M(K)\to E$ is weak$^*$-continuous if and only if there exist a basis $\mu_1,\dots,\mu_n$ of $E$ and functions $g_1,\dots,g_n\in C(K)$ such that
\[
P\nu = \sum_{j=1}^n \Bigl(\int_K g_j\, d\nu\Bigr)\mu_j
\qquad (\nu\in M(K)).
\]
\end{lemma}

\begin{lemma}
\label{lem:ZNm-dense}
The union $\bigcup_{m\in\N} Z_{N,m}$ is dense in $L_1[0,1]$.
\end{lemma}

\begin{theorem}
\label{thm:piecewise-main}
Let $V\subset \ell_1^N$ be a duplication-stable subspace of dimension $n\geqslant 2$, and let $\mathscr V\subset M[0,1]$ be its piecewise-constant copy.

\begin{enumerate}[label=\textup{(\roman*)}]
\item
\[
\lambda(\mathscr V,L_1[0,1]) = \lambda(\mathscr V,M[0,1]) = \lambda(V,\ell_1^N).
\]
\item $\mathscr V$ admits a unique minimal projection from $L_1[0,1]$ onto $\mathscr V$.
\item There is no weak$^*$-continuous minimal projection from $M[0,1]$ onto $\mathscr V$.
\end{enumerate}
\end{theorem}

\begin{proof}
Let $P_V\colon \ell_1^N\to V$ be the unique minimal projection. Define
\[
P_0 := J_N P_V J_N^{-1} A_N\colon L_1[0,1]\to \mathscr V.
\]
Then $P_0$ is a projection onto $\mathscr V$.

Since $A_N$ is contractive and $J_N$ is an isometry,
\[
\|P_0\|\leqslant \|P_V\| = \lambda(V,\ell_1^N).
\]
On the other hand, the restriction of $P_0$ to $Z_N$ equals $J_NP_VJ_N^{-1}$, so
\[
\|P_0\|\geqslant \|J_NP_VJ_N^{-1}\| = \lambda(V,\ell_1^N).
\]
Hence
\[
\|P_0\| = \lambda(V,\ell_1^N).
\]
Because $Z_N\subset L_1[0,1]$, we have
\[
\lambda(\mathscr V,L_1[0,1]) \geqslant \lambda(\mathscr V,Z_N)=\lambda(V,\ell_1^N),
\]
so in fact
\[
\lambda(\mathscr V,L_1[0,1]) = \lambda(V,\ell_1^N).
\]

Now define
\[
\widetilde P_0 := J_N P_V J_N^{-1} R_N\colon M[0,1]\to \mathscr V.
\]
Exactly the same argument yields
\[
\lambda(\mathscr V,M[0,1]) = \lambda(V,\ell_1^N),
\]
proving \textup{(i)}.

We prove uniqueness in $L_1[0,1]$. Fix $m\in\N$. Under the isometry $J_{N,m}$, the subspace $\mathscr V$ corresponds to the duplicated space $D_m(V)\subset \ell_1^{Nm}$.

Let $B_m\colon \ell_1^{Nm}\to D_m(\ell_1^N)$ be the block-averaging projection,
\[
B_m(z_{j,r})_{j,r} = \Bigl(\frac1m\sum_{r=1}^m z_{1,r},\dots,\frac1m\sum_{r=1}^m z_{1,r},\dots,
\frac1m\sum_{r=1}^m z_{N,r},\dots,\frac1m\sum_{r=1}^m z_{N,r}\Bigr).
\]
Then the restriction $P_0|_{Z_{N,m}}$ corresponds, via $J_{N,m}$, to the operator
\[
D_m P_V D_m^{-1} B_m\colon \ell_1^{Nm}\to D_m(V),
\]
which is minimal. Since $V$ is duplication-stable, this minimal projection onto $D_m(V)$ is unique.

Now let $P\colon L_1[0,1]\to \mathscr V$ be any minimal projection. If $P\neq P_0$, choose $f\in L_1[0,1]$ such that $Pf\neq P_0f$. By Lemma~\ref{lem:ZNm-dense}, there exist $m$ and $h\in Z_{N,m}$ with $Ph\neq P_0h$. Because $V$ is duplication-stable,
\[
\lambda(\mathscr V,Z_{N,m}) = \lambda(D_m(V),\ell_1^{Nm}) = \lambda(V,\ell_1^N)=\lambda(\mathscr V,L_1[0,1]).
\]
Hence both restrictions $P|_{Z_{N,m}}$ and $P_0|_{Z_{N,m}}$ have norm equal to $\lambda(\mathscr V,Z_{N,m})$; that is, both are minimal projections from $Z_{N,m}$ onto $\mathscr V$. Via $J_{N,m}$ they correspond to minimal projections from $\ell_1^{Nm}$ onto $D_m(V)$. By uniqueness, they must coincide, a contradiction. Therefore $P=P_0$, proving \textup{(ii)}.

Finally, suppose that $P\colon M[0,1]\to \mathscr V$ is a weak$^*$-continuous minimal projection. By Lemma~\ref{lem:wstar-coded}, there exist a basis $y_1,\dots,y_n$ of $\mathscr V$ and functions $g_1,\dots,g_n\in C[0,1]$ such that
\[
P\mu = \sum_{k=1}^n \Bigl(\int_0^1 g_k\, d\mu\Bigr) y_k
\qquad (\mu\in M[0,1]).
\]
Since $\|P|_{L_1}\|\leqslant \|P\| = \lambda(\mathscr V,M[0,1]) = \lambda(\mathscr V,L_1[0,1])$, the restriction $P|_{L_1}$ is a minimal projection from $L_1[0,1]$ onto $\mathscr V$. By \textup{(ii)},
\[
P|_{L_1} = P_0.
\]

Fix $k\in\{1,\dots,n\}$ and $j\in\{1,\dots,N\}$. We claim that $g_k$ is constant on $I_j$. If not, choose $s,t\in I_j$ with $g_k(s)\neq g_k(t)$. After interchanging $s$ and $t$ if necessary, we may choose $\omega\in\K$ with $|\omega|=1$ such that
\[
\Re(\omega g_k(s)) < \Re(\omega g_k(t)).
\]
By continuity, there exist disjoint subintervals $J,J'\subset I_j$ of the same length such that
\[
\sup_{u\in J} \Re(\omega g_k(u)) < \inf_{v\in J'} \Re(\omega g_k(v)).
\]
Because $J$ and $J'$ have the same length and lie in the same $I_j$,
\[
A_N(\chi_J)=A_N(\chi_{J'}),
\]
hence
\[
P_0(\chi_J)=P_0(\chi_{J'}).
\]
Since $P|_{L_1}=P_0$ and $y_1,\dots,y_n$ is a basis, this implies
\[
\int_J g_\ell(t)\,dt = \int_{J'} g_\ell(t)\,dt
\qquad (1\leqslant \ell\leqslant n).
\]
In particular,
\[
\Re\left(\omega\int_J g_k(t)\,dt\right)
=
\Re\left(\omega\int_{J'} g_k(t)\,dt\right).
\]
On the other hand,
\[
\Re\left(\omega\int_J g_k(t)\,dt\right)
=
\int_J \Re(\omega g_k(t))\,dt
<
\int_{J'} \Re(\omega g_k(t))\,dt
=
\Re\left(\omega\int_{J'} g_k(t)\,dt\right),
\]
a contradiction. Thus each $g_k$ is constant on every $I_j$, and continuity then forces each $g_k$ to be constant on all of $[0,1]$.

Write $g_k\equiv c_k$. Then for $f\in L_1[0,1]$,
\[
P_0f = \sum_{k=1}^n c_k\Bigl(\int_0^1 f(t)\,dt\Bigr) y_k.
\]
Hence the matrix of $P_0|_{\mathscr V}=\Id_{\mathscr V}$ in the basis $y_1,\dots,y_n$ has the form
\[
\bigl(c_k\!\int_0^1 y_\ell(t)\,dt\bigr)_{k,\ell=1}^n,
\]
which has rank at most $1$. Since $n\geqslant 2$, this matrix cannot be the identity. The contradiction shows that no weak$^*$-continuous minimal projection exists, proving \textup{(iii)}.
\end{proof}

\end{document}
